%%%%%%%%%%%%%%%%%%%%%%%%%%%%%%%%%%%%%%%%%
% The Legrand Orange Book
% LaTeX Template
% Version 2.3 (8/8/17)
%
% This template has been downloaded from:
% http://www.LaTeXTemplates.com
%
% Original author:
% Mathias Legrand (legrand.mathias@gmail.com) with modifications by:
% Vel (vel@latextemplates.com)
%
% License:
% CC BY-NC-SA 3.0 (http://creativecommons.org/licenses/by-nc-sa/3.0/)
%
% Compiling this template:
% This template uses biber for its bibliography and makeindex for its index.
% When you first open the template, compile it from the command line with the 
% commands below to make sure your LaTeX distribution is configured correctly:
%
% 1) pdflatex main
% 2) makeindex main.idx -s StyleInd.ist
% 3) biber main
% 4) pdflatex main x 2
%
% After this, when you wish to update the bibliography/index use the appropriate
% command above and make sure to compile with pdflatex several times 
% afterwards to propagate your changes to the document.
%
% This template also uses a number of packages which may need to be
% updated to the newest versions for the template to compile. It is strongly
% recommended you update your LaTeX distribution if you have any
% compilation errors.
%
% Important note:
% Chapter heading images should have a 2:1 width:height ratio,
% e.g. 920px width and 460px height.
%
%	Cities list:
%	New York
%	Beijing
%	London (city)
%	Singapore
%	Dubai
%	Tokyo
%	Los Angeles
%	Frankfurt
% 	Kuala Lumpur
%	Brussel
%	Milan
%	Berlin
%
%%%%%%%%%%%%%%%%%%%%%%%%%%%%%%%%%%%%%%%%%

%----------------------------------------------------------------------------------------
%	PACKAGES AND OTHER DOCUMENT CONFIGURATIONS
%----------------------------------------------------------------------------------------

\documentclass[11pt,fleqn,oneside]{book} % Default font size and left-justified equations

%----------------------------------------------------------------------------------------

\input{structure} % Insert the commands.tex file which contains the majority of the structure behind the template

\begin{document}

%----------------------------------------------------------------------------------------
%	TITLE PAGE
%----------------------------------------------------------------------------------------

\begingroup
\thispagestyle{empty}
\begin{tikzpicture}[remember picture,overlay]
    \node(background) at (current page.center) {\includegraphics[height=\paperheight]{background}};
\end{tikzpicture}
\endgroup

\newpage

% \begingroup
% \section{Abstract}\index{Abstract}
% \thispagestyle{empty}
% \addcontentsline{toc}{chapter}{\textcolor{trolleygrey}{Abstract}}
\usechapterimagefalse % If you don't want to include a chapter image, use this to toggle images off - it can be enabled later with \usechapterimagetrue
\chapter{Abstract}
\textbf{Ripa Exchange è un exchange ibrido-decentralizzato con il focus di abbassare il livello di ingresso per aprire 
nuove piattaforme di exchange ed al contempo offrire ai traders un ambiente sicuro ed efficiente per eseguire operazioni di 
trading giornaliere.}\\

Il Team di Ripa Exchange crede che, indipendentemente dagli sviluppi nel mondo delle criptomonete, è oggigiorno costoso per aprire,
gestire, e trovare fiducia nelle proprie operazioni di cambio valuta non solo per le risorse richieste per avviare
una piattaforma di scambio affidabile ma soprattutto per sviluppare la piattaforma in sè e trovare la liquidità 
necessaria per generare profitto nei primi 3-5 anni di attività.\\

Azione è richiesta ed azione è richiesta adesso. Gli utenti sono frustrati da exchange inaffidabili che chiudono scappando con i 
fondi, vengono hackerati oppure non sostengono il carico giornaliero di una'industria crescente come questa è.
Nonostante l'impegno degli imprenditori di exchange per offrire piattaforme di cambio valuta efficienti, affidabili e 
semplici da utilizzare i fondi necessari per costruire queste piattaforme sono nella fascia di 250-300 mila euro senza includere
il costo del personale per offrire supporto di livello platino agli utenti finali, l'infrastruttura di gestione, e costi giornalieri
dell'attività. Tali costi vengono sostenuti per avere una piattaforma discreta alla quale ci si lega ad unica via con una software
house per sviluppi futuri.\\

É obbiettivo di questo progetto di sviluppare una piattaforma di cambia valuta Open Source, efficiente, affidabile ed 
offrire liquidità\footnote{Grazie alla tecnologia RLSP (Ripa Liquidity Service Provider)} all'exchange così creato dal \textbf{primo}
giorno di attività in modo tale che gli imprenditori in questo settore possono focalizzarsi nel pubblicizzare la propria piattaforma,
offrire supporto tecnico di libello platino ai propri utenti e soddisfare tutte le eterogenee legislazioni nell'industria delle valute virtuali.
Mentre vogliamo che l'esperienza di trading per l'utente finale sia la migliore possibile offrendo un ambiente di trading sicuro ed affidabile.\\
\usechapterimagetrue
% \endgroup

%----------------------------------------------------------------------------------------
%	TABLE OF CONTENTS
%----------------------------------------------------------------------------------------

\chapterimage{chapter_head_1_Beijing.jpg} % Table of contents heading image
\renewcommand*\contentsname{Indice dei Contenuti}
\tableofcontents % Print the table of contents itself

%----------------------------------------------------------------------------------------
%	PART
%----------------------------------------------------------------------------------------


%\part{Part One}

%----------------------------------------------------------------------------------------
%	CHAPTER 2: Introduction
%----------------------------------------------------------------------------------------

\chapterimage{chapter_head_2_London.jpg} % Chapter heading image

\chapter{Introduzione}
Il livello di accesso all'industria delle valute virtuali è piuttosto elevato da un punto di vista tecnico per l'utente medio, 
ed ha un livello di accesso economicamente elevato per l'imprenditore che vuole avviare un'attività in questo campo, valore dato 
dall'acquisto del codice sorgente dell'exchange di criptovalute, dal reclutamento di personale operativo per le operazioni di DevOps,
per reclutare operatori di supporto tecnico, per sottostare alle leggi nazionali ed internazionali in materia di AML/KYC, per avere liquidità
dal primo giorno di avvio delle operazioni di cambio. Con questo progetto, Ripa Exchange, vuole abbassare questo livello di ingresso 
perchè \textbf{gestire un exchange è DIFFICILE} ed il Team Ripa Exchange vuole che il lettore e futuro imprenditore nel campo delle valute
virtuali si concentri/focalizzi su cose veramente importanti non cavilli burocratici che l'industria richiede perchè TU vuoi avviare un'attività
in questo campo e tu necessiti del codice sorgente per avviarla (ed il codice sorgente per avviarla ti deve essere concesso gratuitamente).

Per rafforzare/chiarire questo principio economico il punto è che per avviare un exchange è richiesto un investimento esiguo dal tuo 
venture capital di riferimento e anche con un tale investimento il ritorno dei fondi investiti all'angel investor non sono garantiti
nei primi 5 anni di attività.

Per avviare un servizio professionale di cambio valuta crediamo che il codice sorgente dell'exchange e la liquidità da offrire ai tuoi clienti
dal primo giorno di attività dovrebbero esserti forniti gratuitamente: non deve essere accettabile pagare \euro150.000,00 ad una software house
solo per avere una piattaforma che funzioni e per la quale dovete erogare altri \euro100.000,00-150.000,00 per brandizzarla, personalizzarla
per i vostri scopi legandovi in questo modo ad una compagnia/software house che potrebbe fallire in futuro e che vi farebbe ritrovare
senza poter fare ulteriori modifiche al vostro exchange in quanto non vi hanno mai consegnato il codice sorgente della piattaforma/del prodotto
su cui il vostro business fa affidamento/si basa.

\textbf{Noi crediamo che tutto questo debba essere gratuito}: dobbiamo offrirti la miglior tecnologia nel mercato attuale in modo tale che tu 
possa focalizzarti sulla tua attività mentre noi ci focalizziamo in sviluppare la tecnologia per gestire la tua attività/i tuoi affari
in maniera efficiente, sicura, responsiva e produttiva. Questo è il motivo per cui Ripa Exchange si prefigge lo scopo di sviluppare una rete di exchange
focalizzandosi/utilizzando un'architettura che è \textit{efficiente, sicura, responsiva, compliant e personalizzabile} in modo tale che ogni exchange
nella rete può basarsi su di fondamenta solide mentre è comunque possibile modificare la singola istanza secondo le personalizzazioni 
che la compagnia che gestisce quell'istanza richiede/quel singolo exchange richiede.

Per raggiungere questo scopo abbiamo scelto di sviluppare la tecnologia Ripa Liquidity Service Provider basata su ARK - a blockchain for consumer adoption - 
il cui obbiettivo primario è incrementare l'adozione delle tecnologie blockchain agli utenti finali focalizzandosi su due aree critiche: 
\underline{Una Efficiente e Sicura Tecnologia di Base} e \underline{Servizi Reali per Gente Comune}. L'ecosistema ARK è ancora ad uno stadio
iniziale di sviluppo: la futura implementazione di ARK 2.0 permetterà di eseguire smart contract nativi sulla blockchain, che permetteranno
alla blockchain ARK di competere con la sua concorrente Ethereum da un punto di vista tecnologico. RLSP permetterà a tutti gli exchange della rete Ripa
di condividere la stessa liquidità-lo stesso libro degli ordini in base ai parametri di condivisione della singola istanza dell'exchange Ripa
configurabili nella tua area amministratore.

Il Ripa Founder Team (RFT), come presentato in ripaex.io, opera nel nome del gruppo Ripa. L'RFT è responsaile per l'uso corretto dei finanziamenti
ottenuti nella campagna di scambio token RIPA TEC come spiegato di seguito in questo documento.

Il RFT comprende-impone che i risultati delle fasi di Prevendita e RIPA TEC saranno usati esclusivamente per il finanziamento del progetto 
\emph{Ripa Exchange} come spiegato in questo whitepaper, che saranno resi disponibili per estrazione nella piattaforma apposita resta disponibile
più tardi quest'anno tec.ripaex.io, e che dovranno risultare nella creazione di una entità legale che sarà chiamata \emph{Ripa Exchange}. La creazione di tale 
compagnia è pianificata per il primo quadrimestre del 2019.

Fino a tale data RFT opera invece di \emph{Ripa Exchange, una compagnia nel processo di essere incorporata}.


%------------------------------------------------

%This statement requires citation \cite{article_key}; this one is more specific \cite[162]{book_key}.

%------------------------------------------------

\section{Terminologia}\index{Terminologia}
\begin{description}
	\item[\textsc{Ripa Exchange:}] a FIAT <-> CRYPTO exchange (a cryptocurrency exchange) based on the source code
	of Peatio \cite{peatio}
	\item[\textsc{Ripa Blockchain:}] a DPOS blockchain in which liquidity is exchanged for all the exchanges in the Ripa network
	\item[\textsc{Ripa Token - XPX:}] a cryptographically secure token exchanged on the Ripa blockchain based on the DPOS protocol
	\item[\textsc{RIPA:}] the DPOS financial ecosystem composed of Ripa Exchange and Ripa Blockchain
	\item[\textsc{RIPAEX:}] the name of the project, project website and hosted domain
    \item[\textsc{RLSP:}] Ripa Liquidity Service Provider, a shared orderbook to exchange orders between exchanges in the same Ripa network
	\item[\textsc{ARK:}] a platform for consumer adoption of blockchain technologies \cite{ark}
	\item[\textsc{ACES:}] Ark Contract Execution Services \cite{aces} provides simple protocols and tools for building a robust 
	blockchain service marketplace based on the ARK SmartBridge technology
    \item[\textsc{“,” or “.”:}] The Anglo-Saxon use of decimal points and commas to represent numbers has
been chosen for the purposes of this document: that is to say that a “.” represents a decimal point, and a “,”
distinguishes between multiples of thousands, millions and billions.
    \end{description}

%Lists are useful to present information in a concise and/or ordered way\footnote{Footnote example...}.
\pagebreak
\section{Roadmap}
There are essentially four phases to the RipaEx project:\\

\tcbset{roadmapBox/.style={colback=yellow!10!white,colframe=azure(colorwheel),
equal height group=nobefaf,width=(\linewidth-1pt), height from=4cm to 8cm, nobeforeafter,
	center title,
	valign=top, halign=left}}
\begin{center}
\begin{tcolorbox}[roadmapBox,
	title=\textsc{Funding the project: XPX presale and RIPA TEC (WP2)}]

	This phase recognises the existence of interest in this market development
	from across the World concerning the lowering of the entry level for building a cryptocurrency exchange.
	It aims to make the first comprehensive analysis of this state of the art to form the basis of the later project phases and
	build the first working prototype of a centralized exchange based on Peatio.\\
	\vspace{1cm}
	\centering\textbf{\textsc{Phase ending January 2019}}.
\end{tcolorbox}

\resizebox{0.05\textwidth}{26pt}{$\Downarrow$}

\begin{tcolorbox}[roadmapBox,
	title=\textsc{First exchange opening and development of tools and resources (WP3)}]

	The second phase takes the results of the first 
	and develops from them a set of tools and resources which provide concise and comprehensible guidance to market actors in any
	Country. With the first instance of Ripa Exchange running first contacts with other economical players in the industry can be
	done.\\
	\vspace{1cm}
	\centering\textbf{\textsc{Phase ending June 2019}}.
\end{tcolorbox}

\resizebox{0.05\textwidth}{26pt}{$\Downarrow$}

\begin{tcolorbox}[roadmapBox,
	title=\textsc{Dissemination (WP 7/8) and Project Coordination (WP1)}]

	During the full duration of the project, 
	dissemination activities (WP 7/8) are carried out in which results from the individual work packages are disseminated 
	to relevant target groups including project partners, RipaEx supporters, exchanges managers, banking partners as well 
	as relevant target groups. This phase covers a wide range of dissemination techniques, from printed and
    electronic handbooks to workshops and training sessions, ongoing networks, all having the
	ultimate goal of defining a standard for exchanges communication among public and private entities. 
	An overarching work package is concerned with the management of the project from start to finish, ensuring proper coordination, 
	quality assurance and budgetary control (WP1).
\end{tcolorbox}

\resizebox{0.05\textwidth}{26pt}{$\Downarrow$}

\begin{tcolorbox}[roadmapBox,
	title=\textsc{Development of hybrid-decentralized exchange (WP 4-6)}]

	Using the tools and resources developed in WP3, 
	Work packages 4-6 focus on bringing collected knowledge and tools into practice. The three work packages reflect three major
	focal points (and target groups) within the network of exchange created for establishing successful 
	demonstrations on local scale: incorporations of local Ripa Exchanges (WP4), technical analysis for the 
	Ripa Liquidity Service Provider (WP5), and first MVP of the hybrid decentralized exchange (WP6). The demonstration
    phase forms the heart of the RipaEx action; WP 2 and 3 are focused on providing
	deliverables (e.g. tools) that enable successful and efficient demonstration activities.\\
	\vspace{1cm}
	\centering\textbf{\textsc{Phase ending January 2021}}.
\end{tcolorbox}
\end{center}

\section{RipaEx Partners - RipaEx Governance}\index{RipaEx Partners - RipaEx Governance}
Most of the partners are entrepreneurs in the virtual currency industry, but a research institute 
and Financial Organizations are also represented. The Partners are:
\begin{description}
	\item[Coordinator]: Ripa Exchanges Ltd
	\item[CoBeneficiaries]: 
\end{description}
\subsection{RipaEx Governance}\index{RipaEx Governance}
Governance for the network of exchanges created, development of the source code, owning of the XPX tokens, Ripa Foundation.

\section{Project Summary}
\begin{enumerate}
	\item \textsc{what}: RipaEx is a project to facilitate the uptake of standards to share liquidity between crypto assets marketplaces. 
	The objective of RipaEx is the promotion of shared source code for wallets and exchanges in the virtual currency industry: 
	It is the aim of this reference document to give in-depth information to prospective exchange developers,
	or exchange managers, to enable correct decision-making and to ensure success for their proposed projects. 
	It seeks to analyse the real potential in the Country of application for a network of cryptocurrency exchanges, 
	and its place in the market.
	\item \textsc{what}: Crypto assets are an alternative to centralized assets managed by (country-specific) stock exchanges. Although certain
	stock exchanges gives the possibility to their users to verify and manage the assets they own the verification process
	is not always transparent that is the reason because from 2009 \cite{bitcoin} onwards a new types of (community-verifiable) assets 
	have been implemented to give small, medium and big investors complete transparency in the managing of their investments
	assets.
	\item \textsc{when}: Recent developments at European Union level and worldwide are transforming both how virtual currencies
	are treated and the way ICO (Initial Coin Offering) are legislated. These combined developments 
	have made the use and production of virtual currencies an increasingly favourable prospect. \\
	In October 2015 the European Court of Justice ruled that bitcoin and other cryptocurrencies are exempt from VAT taxation. \\
	In July 2016 the European Commission adopted proposals for legislation to amend the 4th Anti-Money Laundering Directive (4AMLD) that
	will bring virtual currencies exchanged and wallet providers into the EU's anti-money laundering framework \cite{EUAMLCrypto}.\\
	In February 2018 the European Commission launched the EU Blockchain Observatory and Forum \cite{EUBOaF} to highlight key developments 
	of the blockchain technology, promote European actors and reinforce European engagement with multiple stakeholders involved in blockchain activities.
	\item \textsc{why}: However there is still very little regulation performed on ICOs and only United States of America at the moment has undergone
	a legislation defining ICO tokens as securities. \cite{SECICO}
	\item \textsc{who}: The results indicate that medium tech savvy from 18 to 45 is the average user of
	virtual currencies although the corporate finance companies are also starting to put 
	virtual currencies schemes inside their portfolio especially since the presentation 
	of the bitcoin futures contract from CME Group Inc. in the stock exchange of Chicago
	last 18th of December 2017.
	\item \textsc{how much}: Total virtual currencies market capitalization has been estimated around 317 B USD\footnote{Coinmarketcap data April 2018}
	and is predicted to grow to 5,000.00 B USD in the next ten years span \cite{cryptoMCTenYears}.
	\item \textsc{where}: Local authorities are working with National Governments to make sure local exchangers
	in the national territory are complying with national and international AML/KYC regulations.
	Venture capitals and Angel Investors are starting to release financing solutions to start-ups 
	in the Fintech industry all over the world from America to Asia passing through Europe and some 
	Countries are starting state-owned cryptocurrencies schemes to test the exchange of goods \& services
	on those (distributed ledger) technologies \cite{petro}.
	\item \textsc{how much}: The average cost for starting your own crypto asses marketplace is around \$ 150,000.00 only for a running instance of your exchange platform:
	to that you need to add costs to customize the platform before launch and in the future, advertising your new business, running costs for servers,
	network operators, support center operators and legal department to comply with your State of incorporation AML/KYC legislations and general company laws.\\
	\textbf{That is the reason because we think owning the source code of your exchange software is the best way to run a business in this industry}.
	\item \textsc{how}: The main problems encountered in opening a FIAT <-> CRYPTO marketplace is to find trusted
	banking partners to comply with the many different AML/KYC rule and procedures to exchange virtual currencies
	to FIAT currencies.
	\item \textsc{what}: Classical types of exchanges operations are: 
		\begin{enumerate}[label*=\arabic*.]
			\item \textbf{one-way exchanges}: in which a centralized application has all the liquidity to offer to its potential users
			\item \textbf{two-way exchanges}: in which a centralized or decentralized platform match the selling requests with the buying requests
			of its users
		\end{enumerate}
		On this a sub-classification is also necessary:
		\begin{enumerate}[label*=\arabic*.]
			\item \textbf{FIAT <-> CRYPTO exchanges}: in which exchanges operations are performed between FIAT\footnote{Traditional central banks owned currencies like EUR, USD, GBP, JPY, others...} 
			currencies and virtual currencies
			\item \textbf{CRYPTO <-> CRYPTO exchanges}: in which exchanges operations are performed only between virtual currencies
		\end{enumerate}
	You can build a matrix based on the four configurations above to build the exchange operation platform of your needs.
	\item \textsc{with what}: The specifications to look when choosing for an exchange platform to run are:
		\begin{enumerate}[label*=\arabic*.]
			\item \textbf{code}: Open Source, Closed Source or hybrid solution
			\item \textbf{modularization}: separation between exchange engine (orders matching engine), UI and user registry
			\item \textbf{UI responsiveness}
			\item \textbf{compliance} with current industry standards
			\item \textbf{customization} of the exchange engine, trading currencies, UI and other aspects of the crypto asset marketplace platform...
			\item \textbf{security} of the funds: saving in cold wallets and hot wallets configurable
			\item \textbf{transparency} of the funds: proof of solvency of the exchange
			\item \textbf{Multi-Accounts trading}: easy to configure new virtual currency protocols
			\item \textbf{Multi-Accounts users}: possibility to interact with user accounts from Google, Facebook, Twitter 
			to login into the platform and FIDO Alliance security standards for personal credentials.
		\end{enumerate}	
	\textbf{Those are not only technical decisions to be made but also economical} especially the owning of the source code of your crypto asset
	marketplace platform is fundamental to make future customization of your exchange in an independent way compared to rely on a single
	software house that makes the customizations for you.
	\item \textsc{how}: Options for finding users for your exchanges operations are: targeted marketing campaigns, innovative features in the industry,
	fee level based on trading quantities, bonuses for first registration and trading quantities, affiliate marketing for paying users to take their friends
	to your exchange.
	\item \textsc{how}: For setting-up a crypto asset marketplace a project must take into account the following legislation:
		\begin{enumerate}[label*=\arabic*.]
			\item \textbf{AML/KYC}: \textit{Fourth Anti-Money Laundering Directive} if business set up in the European Union \cite{4AMLD} or the AML/KYC reference implementation
			to your crypto asset marketplace Country of incorporation (as an example \textit{Intelligence Reform \& Terrorism Prevention Act of 2004}
			written by FinCEN in the United States of America).\\
			International recommendations for undergoing AML/CFT verifications are given by the Financial Action Task Force on Money Laundering \cite{FATF}.
			\item \textbf{Payment Licence}: By far the biggest and most arduous task with regards to legitimising the 
			FIAT <-> CRYPTO exchanges operations is obtaining a \textit{PSD Licence} \cite{PSD}. 
			The PSD licence follows Council Directive 2007/64/EC and is applied in each country via its own national laws. 
			Costs of an IPPC licence can vary between \euro XXXX and \euro XXXX, depending on the size of operation.
		\end{enumerate}
	\item \textsc{who}: The nature of the business under consideration by the Ripa Exchange project (small scale,
	localised FIAT <-> CRYPTO exchanges operations), means that each enterprise likely to have 7 or 8 staff: N.2 developers, 
	N.1 network/security operator, N.1 administrative, N.2 client support operators, N.1 legal and tax advisor. \\
	The turnover of such an enterprise however, because of the high value of the end product, is likely to be more than 
	\euro 350,000 a year and could be several times higher. A business of
	this scale lends itself to the following possible company structures: A simple partnership;
	A limited company; A non-profit company or social enterprise; A worker co-operative.
	Financial Agencies are potential key actors, but the type of business they can set up will
	depend on their legal status which does vary from country to country.
	\item \textsc{how}: Potential sources of funds for a small-medium sized crypto asset marketplaces are: Bank Loans; Low
	Interest Loan Schemes; Commercial Credit; Equity financing; Business Angels venture
	capital. Having a robust Business Plan and financial guarantees are essential elements
	for securing funding. The European Investment Fund (EIF) of the EIB, offers support in
	the form of guarantees for SMEs.
	\item \textsc{why}: The arguments for crypto asset marketplaces are for financial freedom, decentralizing of the value-transferring operations, and
	owning for real your money. 
	There are other benefits, well documented, such as faster payments, long term gain based on deflationary economy and prediction of 
	Great Depressions like the one that hit the global economy in 2008. But
	above all, virtual currencies are the only direct competitor to centralized value-transferring operations done by central banks.

	\item \textsc{why}: There is consensus in the literature that the use of virtual currencies in place of fiat currencies will result in 
	higher financial freedom especially as they fit into the Austrian school of economy \cite{austrianTheory} 
	(TODO: add more on Austrian economics school)

	\item \textsc{why}: Benefits of virtual currencies schemes (TODO: put some numbers)

	\item \textsc{how}: Securing assets on the blockchains means basically performing three operations
		\begin{enumerate}[label*=\arabic*.]
			\item \textbf{Generating a random private key}
			\item \textbf{Converting the private key generated in (1) into a public key}: a common protocol making this conversion
			in the virtual currencies industry is the ECDSA curve algorithm
			\item \textbf{Converting the public key generated in (2) into a virtual currency address}: common protocols for making this conversion
			are hash functions SHA-256, Base58 encoding, Base32 encoding
		\end{enumerate}
	At this point any value sent to the virtual currency address generated in (3) is secured on the blockchain of choice and accessible
	only from the owner of the relative private key generated in (1).
	\item \textsc{where}: The two critical factors affecting the cryptocurrency industry are banks concurrence and State banning.
	Although a harmonisation throughout Europe would be beneficial to development of the industry both in terms of 
	taxation and warranty approvals, this is currently not the case. Each country has its specific legislation and tax
	regime for all exchanges operations involving FIAT money, and State banning is going to completely liberalization 
	of this activities like European Union to complete banning and imprisonment of operators in this industry like Bangladesh
	\cite{bitcoinLegality}. 

	\item \textsc{where}: The Asiatic countries of South Korea, China and Japan are the leader in the field of cryptocurrencies for number of 
	transactions for over 9 years with a proactive approach and favourable tax regime. At the beginning of 2017 in Japan bitcoin
	has been declared legal tender but China has recently declared illegal token sale and exchanges and local cryptocurrencies
	marketplaces are closing down.
	\item \textsc{where}: Any assessment of your local market should include: number of potential users to reach, 
	type of exchange to incorporate (FIAT <-> CRYTPO or CRYPTO <-> CRYPTO), type of virtual currencies protocol 
	to integrate (POW, DPOS, Masternodes, others...), types of services to offer (exchange only, advanced trading tools,
	payment processor, others...), if FIAT <-> CRYPTO exchange number of FIAT payments processors to accept (PayPal, OKPay,
	MoneyPolo, others...), number of others exchanges in your region.

	\item \textsc{who}: there are a number of options for dealing with Warranty/Customer protection issues: 
	creating consumer pressure by making clear to the end users that the possession of the private keys of their
	virtual currency addresses make \textbf{liable} for any loss of the private keys meaning nobody can help
	them recovering their funds if the their private keys are lost. Creating consumer pressure to not leave funds
	on exchanges ("\textit{Be Your own Bank!!}"), making them choose the licensed exchanges in the market. 
	\item \textsc{who}: While it is very expensive to insure money exchanges operations and money transmitting operations, 
	examples of customer protections in the industry are: Kraken platform which is offering Mt. Gox users partial refund
	of their lost, NEO community giving refund to the users involved in the BitGrail hacking, Ethereum supporters
	giving The DAO investors partial refunds, other hacking cases...

	\item \textsc{Recommendations for all/law compliance}: if you intend to incorporate a FIAT <-> CRYPTO exchange you should
	focus from the first instance on law compliance by studying the AML/KYC laws of the country of incorporation and
	finding bank partners to work with. Local financial Authority can help to comply with rules \& regulations and 
	local cryptocurrencies foundations can help you to tune your exchanges operations to perform targeted
	operations based on the customers interests in the country of incorporation.
	Promote cryptocurrency-friendly users in the area of interest.
\end{enumerate}


%----------------------------------------------------------------------------------------
%	CHAPTER 3: Ripa Exchange
%----------------------------------------------------------------------------------------

\chapterimage{chapter_head_3_Singapore.jpg} % Chapter heading image

\chapter{The Ripa Exchange}
Ripa Exchange is an crypto asset (fiat money or cryptocurrency or something) marketplace, following the latest 
industry standards and resting in the principles of "\textit{open source, secure and efficient}". Ripa Exchange
aims to serve a platform for crypto-currency enthusiasts by providing a safe, secure, UI responsive, customizable 
and easy to use exchange that embraces open source principles and public trust.\\\\
Ripa Exchange is implemented with the Rails framework and other cutting-edge technology and will be migrated to an
hybrid-decentralized exchange where all exchanges in the Ripa network will share liquidity thanks to the RLSP technology.

\section{Mission}
\begin{quotation}
	``\textit{Our mission is to build the world best open source crypto asset marketplace with a high performance trading engine 
	and safety which can be trusted and enjoyed by users. Additionally we want to move the crypto currency exchange technology 
	forward by providing support and add new features. We are helping people to build easy their own exchange around the world.}''
\end{quotation}

Help is greatly appreciated, feel free to submit pull-requests or open issues.

\section{Features}
A free, transparent and internationalized open source crypto currency exchange.\\

\tcbset{featureBox/.style={colback=yellow!10!white,colframe=white!20!black,
equal height group=nobefaf,width=(\linewidth-1pt)/3, height=4.5cm, nobeforeafter,
	center title,
	valign=top, halign=left}}
\begin{tcolorbox}[featureBox,
	title=\textsc{Open Source}]

	\small	All source code are fully released under the terms of the MIT License.\\\vspace{5mm}
	\tiny Ripa Exchange is a customizable cryptocurrency exchange solution architecture enables easy connection to KYC/AML, 
	authentication, ETL/reporting, and other services.
\end{tcolorbox}
\begin{tcolorbox}[featureBox,
	title=\textsc{Compliant}]

	\small	International KYC/AML standards.\\\vspace{5mm}
	\tiny Ripa Exchange KYC efficiently submits and exchanges KYC information 
	to meet the banking supervisory standards and comply with Customer Due Diligence (CDD) requirements.
\end{tcolorbox}
\begin{tcolorbox}[featureBox,
	title=\textsc{Transparent \& Configurable}]

	\small	Customize in your own way\\\vspace{5mm}
	\tiny Major functions have been embedded in the source code – neat registration and log-in interface, 
	personalized deposit and withdraw procedure, best match of bid and ask, etc. These functions are comprehensive 
	and are ready to use with no extra work needed. 
\end{tcolorbox}

\begin{tcolorbox}[featureBox,
	title=\textsc{Internationalization}]

	\small	All users are able to view Ripa Exchange in a language to their best convenience.\\\vspace{5mm}
	\tiny Supporting many common languages, Ripa Exchange makes it easy for users to operate in their mother tongue. 
	You are encouraged to contribute to our language variety. Users will benefit from your efforts.
\end{tcolorbox}
\begin{tcolorbox}[featureBox,
	title=\textsc{Proof of Solvency}]

	\small	Easy deployable PoS.\\\vspace{5mm}
	\tiny Ripa Exchange Proof of Solvency (PoS) allows users to verify 
	the solvency of the Ripa Exchange based cryptocurrency exchange without compromising user privacy.
\end{tcolorbox}
\begin{tcolorbox}[featureBox,
	title=\textsc{Multi-Accounts Trading}]

	\small	Easy currency configuration.\\\vspace{5mm}
	\tiny Ripa Exchange allows to create multiple accounts and trading in multiple currencies. 
	Ripa Exchange makes it is easy to trade different currencies.
\end{tcolorbox}

\begin{tcolorbox}[featureBox,
	title=\textsc{Multi-Accounts Users}]

	\small	Easy account configuration.\\\vspace{5mm}
	\tiny Ripa Exchange allows to create multiple login accounts Google, Facebook, Twitter and FIDO Alliance
	login standards to secure your account.
\end{tcolorbox}
\begin{tcolorbox}[featureBox,
	title=\textsc{Enterprise Exchange}]

	\small	Start small, grow big.\\\vspace{5mm}
	\tiny Ripa Exchange enterprise exchange features include a high-performance matching engine, 
	scalable distributed worker threads, and SMS 2-factor authentication.
\end{tcolorbox}
\begin{tcolorbox}[featureBox,
	title=\textsc{Functional \& Intuitive}]

	\small	For the new trader, for the experienced trader.\\\vspace{5mm}
	\tiny Clean, user friendly registration and login interface. 
	Personalized deposit and withdraw procedure and a built-in proof-of-solvency audit.
\end{tcolorbox}

\section{Functional Analysis}
\begin{enumerate}
	\item \textsc{Why Redis-RabbitMQ is needed}:
	\item \textsc{Why NodeJS}: 
	\item \textsc{Why Pusher}:
	\item \textsc{ER MySQL}:
	\item \textsc{Ruby folders hierarchy}:
\end{enumerate}

\section{Technology Stack}
\begin{description}
	\item[\textsc{Ruby on Rails}]:
	\item[\textsc{MySQL}]: 
	\item[\textsc{Redis}]:
	\item[\textsc{RabbitMQ}]:
	\item[\textsc{NodeJS}]:
	\item[\textsc{Pusher}]:
\end{description}

\section{User Interface}
Ripa Exchange user interface is based on Peatio user interface a UI responsive user interface built in Ruby on Rails and 
completely separated from the exchange order engine.
Design for a customised UI are in place to offer to Ripa Exchange users the best experience on all devices: here you can find
some screenshot of the current exchange interface design.

\subsection{End-User Interface}
Following some screenshot of the end-user Ripa Exchange-Peatio interface: 
\textcolor{darkred}{those screenshots are a work in progress and they may not represents the user interface of the final product}\\
\begin{center}
	\includegraphics[width=\textwidth*2/3,height=\textheight,keepaspectratio]{RE_TradingUIC}
	\captionof{figure}{Ripa Exchange trading UI}
	\includegraphics[width=\textwidth*2/3,height=\textheight,keepaspectratio]{RE_depositBTCC}
	\includegraphics[width=\textwidth*2/3,height=\textheight,keepaspectratio]{RE_historyC}
	\includegraphics[width=\textwidth*2/3,height=\textheight,keepaspectratio]{RE_solvencyC}
	\captionof{figure}{Ripa Exchange deposit/whitdraw, order history and solvency screen}
\end{center}

\subsection{Admin Interface}
Following some screenshot of the administrative console of Ripa Exchange-Peatio interface: 
\textcolor{darkred}{those screenshots are a work in progress and they may not rapresents the user interface of the final product}\\
\begin{center}
	\includegraphics[width=\textwidth*2/3,height=\textheight,keepaspectratio]{RE_adminDashboardC}
	\includegraphics[width=\textwidth*2/3,height=\textheight,keepaspectratio]{RE_adminDepositsC}
	\includegraphics[width=\textwidth*2/3,height=\textheight,keepaspectratio]{RE_adminWithdrawsC}
	\captionof{figure}{Ripa Exchange Admin console dashboard, deposit/whitdraw screen}
	\includegraphics[width=\textwidth*2/3,height=\textheight,keepaspectratio]{RE_adminUserInfo1C}
	\includegraphics[width=\textwidth*2/3,height=\textheight,keepaspectratio]{RE_adminVerifyAccountC}
	\captionof{figure}{Ripa Exchange Admin user info and verify account screen}
\end{center}

\section{Features at Launch}
\begin{description}
	\item[\textsc{CRYPTO <-> CRYPTO}]:
	\item[\textsc{OAuth}]: Facebook, Google, Twitter
	\item[\textsc{FIDO}]: login with FIDO Alliance standards
	\item[\textsc{Protocols}]: POW, DPOS, Masternode, tether, ERC20
	\item[\textsc{Currencies}]: BTC, ETH, DOGE, BCH, TUSD, ARK, LISK, SHIFT, RISE, KAPU, OXY, RIPA, promising ERC20 tokens
	\item[\textsc{Main Markets}]: BTC, ETH, ARK
	\item[\textsc{Orders Type}]: market, limit
\end{description}

\subsection{Future Features}
\begin{description}
	\item[\textsc{E-Wallets}]: OKPay, NETELLER, MoneyPolo, others...
	\item[\textsc{Advanced Trading Features}]: margin trading, stop loss and take profit, 
	\item[\textsc{FIAT <-> CRYPTO}]
	\item[\textsc{Other tools}]: VISA/MasterCard, merchants tools, P2P Lending 
\end{description}


\section{Towards a Decentralized Exchange...}
Ripa Exchange is a centralized exchange which will be converted into an hybrid-decentralized exchange to create a network of exchanges
that share the same liquidity between each of them so you can offer liquidity from day 1 of your exchanges operations.\\

To offer the PROs of a centralizide exchange with all the PROs of a decentralized exchange without the CONs of both of them, 
we need to convert first into an hybrid-decentralized exchange and during phase 3 of our project (WP4-6) we will make all the 
functional and technical analysis required to make the next step on solid ground.

%----------------------------------------------------------------------------------------
%	PART
%----------------------------------------------------------------------------------------

%\part{Part Two}

%----------------------------------------------------------------------------------------
%	CHAPTER 4: Ripa Blockchain
%----------------------------------------------------------------------------------------

\chapterimage{chapter_head_4_Dubai.jpg} % Chapter heading image

\chapter{La Blockchain Ripa}

\section{Table}\index{Table}

\begin{table}[h]
	\centering
	\begin{tabular}{l l l}
		\toprule
		\textbf{Treatments} & \textbf{Response 1} & \textbf{Response 2} \\
		\midrule
		Treatment 1         & 0.0003262           & 0.562               \\
		Treatment 2         & 0.0015681           & 0.910               \\
		Treatment 3         & 0.0009271           & 0.296               \\
		\bottomrule
	\end{tabular}
	\caption{Table caption}
\end{table}

%------------------------------------------------

\section{Figure}\index{Figure}

\begin{figure}[h]
	\centering\includegraphics[scale=0.5]{placeholder}
	\caption{Figure caption}
\end{figure}


%----------------------------------------------------------------------------------------
%	CHAPTER 5: Token Sale
%----------------------------------------------------------------------------------------

\chapterimage{chapter_head_5_Tokyo.jpg} % Chapter heading image

\chapter{Lo Scambio di Token XPX}

\section{Table}\index{Table}

\begin{table}[h]
	\centering
	\begin{tabular}{l l l}
		\toprule
		\textbf{Treatments} & \textbf{Response 1} & \textbf{Response 2} \\
		\midrule
		Treatment 1         & 0.0003262           & 0.562               \\
		Treatment 2         & 0.0015681           & 0.910               \\
		Treatment 3         & 0.0009271           & 0.296               \\
		\bottomrule
	\end{tabular}
	\caption{Table caption}
\end{table}

%------------------------------------------------

\section{Figure}\index{Figure}

\begin{figure}[h]
	\centering\includegraphics[scale=0.5]{placeholder}
	\caption{Figure caption}
\end{figure}

%----------------------------------------------------------------------------------------
%	CHAPTER 6: Conclusion
%----------------------------------------------------------------------------------------

\chapterimage{chapter_head_6_Frankfurt.jpg} % Chapter heading image

\chapter{Team \& Conclusion}


%----------------------------------------------------------------------------------------
%	CHAPTER 7: Legal
%----------------------------------------------------------------------------------------

\chapterimage{chapter_head_7_LosAngeles.jpg} % Chapter heading image

\chapter{Termini e Condizioni}

%----------------------------------------------------------------------------------------
%	BIBLIOGRAPHY
%----------------------------------------------------------------------------------------

\chapterimage{chapter_head_8_KualaLumpur.jpg} % Chapter heading image

\addcontentsline{toc}{chapter}{\textcolor{trolleygrey}{Riferimenti}}
\nocite{*}
\printbibliography[title={Riferimenti}]
%------------------------------------------------


%----------------------------------------------------------------------------------------
%	COPYRIGHT PAGE
%----------------------------------------------------------------------------------------

\newpage
~\vfill
\thispagestyle{empty}

\noindent Copyright \copyright\ 2018 Ripa Founder Team\\ % Copyright notice

\noindent \textsc{Pubblicato da Ripa Founder Team}\\ % Publisher

\noindent \textsc{www.ripaex.io}\\ % URL

\noindent Materiale protetto dalla Licenza MIT (di seguito la ``Licenza''). Non dovrai usare questo file eccetto in conformità con la Licenza. Puoi ottenere una copia della Licenza a \url{https://opensource.org/licenses/MIT}. A meno che non sia richiesto dalla lagislazione applicata o concordato via scritto, il software distribuito con la Licenza è distribuito \textsc{``così com'è'' senza garanzie di alcun tipo}, esplicite o implicite. Fate riferimento alla Licenza per il linguaggio specifico che governa i permessi e le limitazione indicati dalla Licenza.\\ 

\noindent \textit{Prima pubblicazione, Marzo 2018} % Printing/edition date

\end{document}
